% TOP_PROBLEM: {"id": 1206, "title": "Chern's Affine Conjecture", "category_id": 6, "category": {"id": 6, "name": "geometry", "display_name": "Geometry", "description": "Euclidean and non-Euclidean geometry, geometric structures.", "slug": "geometry", "order_index": 6, "created_at": "2026-07-31T15:26:25.671Z"}, "status": "open"}
% FIRST_POSTED: 2026-09-27
% !TeX program = lualatex
% ATTEMPT_STATUS: unresolved
% Model/process: GPT_6_Astra_Ultra; autonomous mathematical attempt,
% primary-source audit, exact finite diagnostics, and adversarial review.
% Prepared 27 September 2026; no general solution or novelty claimed.
\documentclass[11pt]{article}
\usepackage[margin=25mm]{geometry}
\usepackage{fontspec}
\setmainfont{Latin Modern Roman}
\usepackage{amsmath,amssymb,mathtools}
\usepackage{unicode-math}
\setmathfont{Latin Modern Math}
\usepackage{xurl,hyperref}
\hypersetup{colorlinks=true,urlcolor=blue}
\setcounter{secnumdepth}{0}
\setlength{\parindent}{0pt}
\setlength{\parskip}{5pt}
\setlength{\emergencystretch}{3em}
\begin{document}
\section*{\texorpdfstring{Chern's Affine Conjecture}{Chern's Affine Conjecture}}\label{1206}
\subsection{Definitions and mathematical statement}
An affine structure on an $n$-manifold is an atlas whose transition maps are restrictions of affine transformations of ${\mathbb{R}}^n$. Equivalently, its tangent bundle has a connection $\nabla$ with zero curvature and zero torsion. A closed manifold is compact without boundary. Its Euler characteristic is
\[
 \chi(M)=\sum_{i=0}^n(-1)^i\dim_{{\mathbb{Q}}} H_i(M;{\mathbb{Q}}).
\]
Chern's usual conjecture, in positive dimension, is that every closed affine manifold satisfies $\chi(M)=0$. Completeness of the affine connection is not assumed; this distinguishes the question from compact complete affine manifolds.

The positive-dimensional convention needs to be explicit: if the catalogue's unqualified phrase every closed affine manifold includes dimension zero, a single point has a flat torsion-free tangent connection and Euler characteristic one. That degenerate literal extension is not the standard conjecture.
\par\smallskip\textit{Formulation and concept references:} \href{https://link.springer.com/article/10.1007/s41884-026-00193-8}{[S1]}.

\subsection{Short English statement}
Must a closed positive-dimensional manifold admitting affine coordinate changes have Euler characteristic zero, even without completeness?
\subsection{Research attempt}

\paragraph{Scope and status audit (27 September 2026).}
The manifold is smooth, closed, and of positive dimension. It is enough to
treat connected manifolds, since Euler characteristic adds over components.
Passing to a finite cover of degree $q$ multiplies Euler characteristic by
$q$ and pulls back the affine structure. Thus an orientation double cover
reduces the problem to the oriented case. No developing map is assumed
injective or surjective, and no completeness hypothesis is inserted.

The supplied 2026 source [S1, Remark 2.4] explicitly states the conjecture.
Klingler's published theorem [E1] proves it under the additional hypothesis
of a parallel volume form. That theorem is acknowledged, not rederived
here. Status searches must distinguish a proof claim from an accepted
result: Ge's arXiv record [E2] is explicitly withdrawn and reports errors,
including an interchange of limits. Its unchanged abstract cannot be used
as a theorem. Cocos's arXiv record [E3] contains a general proof claim;
the present note has not audited that argument and does not use it. These
checks justify retaining the catalogue's unresolved/disputed-claim status,
not announcing a settlement from a search result.

The main route pursued here repairs the tempting but invalid instruction
``put the flat curvature into the Euler Pfaffian.'' A compatible metric
connection is constructed exactly, its residual curvature is computed,
and a sufficient analytic vanishing criterion is derived. The construction
also supplies a local counterexample to pointwise Euler-density vanishing.
Alongside it, elementary holonomy subcases and the two-dimensional
obstruction are checked to delimit where a counterexample could occur.

\paragraph{Elementary reductions which do not need completeness.}
For an oriented closed manifold of odd dimension $n$, rational Poincare
duality gives $b_i=b_{n-i}$. In the Euler sum, the terms with indices $i$
and $n-i$ have opposite signs, and there is no middle index. Hence
$\chi(M)=0$. Applying this on an orientation double cover proves the same
for nonorientable odd-dimensional closed manifolds. The remaining issue
is therefore even dimension.

\textbf{Lemma 498.1 (invariant line subcase).}
If the linear holonomy of a closed affine manifold preserves a
one-dimensional subspace, then its Euler characteristic is zero.

\emph{Proof.}
Parallel transport of the invariant subspace defines a rank-one subbundle
$L\subset TM$: path dependence is precisely the holonomy action, which
preserves the chosen line. Pass, if necessary, to the double cover which
orients $L$. An oriented real line bundle is trivial: choose a bundle
metric and take the positive unit vector in each fiber. This is a smooth
nowhere-zero section of $L$, hence of the tangent bundle of the cover.
The Poincare--Hopf theorem gives Euler characteristic zero there. Dividing
by the finite covering degree proves the assertion for $M$.
\hfill$\square$

The parallel transport need not fix a vector, and may rescale it; the
proof intentionally uses a line subbundle rather than incorrectly claiming
that a parallel vector field exists. More generally, a finite holonomy
orbit of lines becomes an invariant line after passing to the finite cover
corresponding to its stabilizer. The same argument applies.

If the closure of linear holonomy is compact, average a positive inner
product over this compact group and extend it by parallel transport.
Holonomy invariance makes the resulting Riemannian metric globally
well-defined and parallel. The flat connection is now a metric connection,
so its Euler form vanishes, giving $\chi(M)=0$ by Chern--Gauss--Bonnet.
Here the averaging is over a compact group with finite Haar measure; there
is no analogous averaging over an arbitrary noncompact linear holonomy
group. The compact case is one place where the tempting curvature argument
does become valid after its missing hypothesis is verified.

\paragraph{Surface check using the precise flat-bundle inequality.}
For a closed oriented surface $\Sigma_g$ with $g\ge1$, Milnor's Theorem 1
[E4] implies that a flat oriented real two-plane bundle $E$ satisfies
\[
 |\langle e(E),[\Sigma_g]\rangle|<g.
\]
For $E=T\Sigma_g$ its Euler number is $2-2g$. If $g\ge2$, then
$|2-2g|=2g-2\ge g$, a contradiction. If $g=0$, a flat bundle on the simply
connected sphere has a global parallel frame by path-independent parallel
transport, whereas the tangent bundle has Euler number two. Thus $g=1$
is the only oriented possibility and has Euler characteristic zero. The
orientation cover handles every nonorientable closed surface.

This is an application of a cited theorem, not an independent proof of
Milnor's inequality. The indexed primary article supplied its exact strict
bound and corollary; direct retrieval from its archive encountered an access
barrier, recorded in [E4]. The inequality is enough for the surface
conclusion even for a flat tangent connection with torsion. In higher
dimensions that extra strength cannot simply be assumed.

\paragraph{The metric correction and its surviving curvature.}
Let $\nabla$ be a flat connection on $TM$, and choose any smooth positive
definite bundle metric $g$. The following identities use flatness but not
yet torsion-freeness. Define an endomorphism-valued one-form
\[
 B=g^{-1}\nabla g,\qquad A=\tfrac12 B,
 \qquad D=\nabla+A.
\]
The notation means that
$g(B_XY,Z)=(\nabla_Xg)(Y,Z)$, so each $B_X$ is self-adjoint with respect
to $g$. This definition is intrinsic; the matrix formulas below are
computed in a local $\nabla$-parallel frame.

\textbf{Lemma 498.2 (exact correction identity).}
The connection $D$ is compatible with $g$, and its curvature is
\begin{equation}\label{ra498:curvature}
 R^D=-A\wedge A=-\tfrac14B\wedge B,
 \qquad
 R^D(X,Y)=-[A_X,A_Y].
\end{equation}

\emph{Proof.}
Since $\nabla_Xg=gB_X$ in matrix notation and $B_X$ is $g$-self-adjoint,
\[
 (D_Xg)(Y,Z)=(\nabla_Xg)(Y,Z)
              -g(A_XY,Z)-g(Y,A_XZ)=0.
\]
In a parallel frame $B=g^{-1}\mathrm dg$. Differentiating the inverse
matrix gives $\mathrm d(g^{-1})=-g^{-1}(\mathrm dg)g^{-1}$, hence
\[
 \mathrm dB=-B\wedge B,
 \qquad \mathrm dA=-2A\wedge A.
\]
The original connection matrix is zero in this frame, and therefore
$R^D=\mathrm dA+A\wedge A=-A\wedge A$. Because both sides transform
tensorially, the identity is global. Evaluating the wedge product on two
vectors gives the displayed commutator formula.\hfill$\square$

When $n=2r$, Chern--Weil theory applied to the \emph{metric} connection $D$
represents $e(TM)$ by its Euler form
\begin{equation}\label{ra498:euler}
 e(D)=\frac{1}{(2\pi)^r}\operatorname{Pf}(\Omega_D),
 \qquad \chi(M)=\int_M e(D),
\end{equation}
where $\Omega_D$ is the skew-symmetric curvature matrix in an oriented
$g$-orthonormal frame, with a consistent curvature sign convention.
The Pfaffian is not being applied to an arbitrary matrix in an arbitrary
parallel frame. Formula \eqref{ra498:curvature} shows exactly why flatness
of $\nabla$ does not imply flatness of the compatible connection: the
endomorphisms measuring the metric variation need not commute.
The Euler-form construction and the nonmetric-connection distinction are
discussed in [E5, Appendix C].

\paragraph{A local stress test, even with a Hessian metric.}
On a disk about the origin in $\mathbb R^2$, use the ordinary flat
torsion-free connection and the potential
\[
 f(x,y)=\tfrac12(x^2+y^2)+\tfrac16x^3-\tfrac12xy^2.
\]
Its Hessian is
\[
 g=\begin{pmatrix}1+x&-y\\-y&1-x\end{pmatrix}
  =I+xS+yT,\quad
 S=\begin{pmatrix}1&0\\0&-1\end{pmatrix},\quad
 T=\begin{pmatrix}0&-1\\-1&0\end{pmatrix}.
\]
The eigenvalues are $1\pm\sqrt{x^2+y^2}$, so $g$ is positive definite
on the unit disk. At the origin the coordinate frame is orthonormal and
\[
 A_{\partial_x}=S/2,\qquad A_{\partial_y}=T/2,
 \qquad
 R^D(\partial_x,\partial_y)
 =-\tfrac14[S,T]
 =\begin{pmatrix}0&1/2\\-1/2&0\end{pmatrix}.
\]
Thus the degree-two Euler density at the origin is nonzero (its magnitude
is $1/(4\pi)$ in the coordinate area form). In this example the third
derivatives of $f$ are totally symmetric, so
$A_XY=A_YX$ for coordinate vectors. Hence $D$ is also torsion-free and is
the Levi--Civita connection of $g$. The nonzero density is not an artifact
of allowing torsion in the corrected connection.

This is not a counterexample to the conjecture: the disk is not closed.
It is a counterexample to the stronger local claim that any metric Euler
density must vanish on an affine chart. By inserting a cutoff of a small
multiple of $xS+yT$ in a coordinate disk on a flat two-torus, one obtains
a globally positive metric with the same nonzero commutator near the disk
center. The torus still has Euler characteristic zero. Its Euler integral
therefore cancels across the manifold. A successful general argument must
control a global integral or cohomology class, not merely use the existence
of locally flat coordinates.

\paragraph{A quantitative sufficient criterion and its scaling obstruction.}
The exact formula yields an analytic route with a clearly measurable
target. Let $|A|_g$ denote the Hilbert--Schmidt norm on
$T^*M\otimes\operatorname{End}(TM)$, using $g$ on all factors. At each
point choose a $g$-orthonormal basis. For basis vectors $e_i,e_j$,
\[
 |R^D(e_i,e_j)|\le2|A_{e_i}|\,|A_{e_j}|\le2|A|_g^2.
\]
The coefficient of a Pfaffian is a finite sum, depending only on $r$, of
products of $r$ curvature components. Consequently a dimensional constant
$C_r<\infty$ satisfies
\begin{equation}\label{ra498:energy}
 |\chi(M)|\le\int_M|e(D)|
 \le C_r\int_M|A|_g^{2r}\,\mathrm d\operatorname{vol}_g.
\end{equation}
This estimate follows just from the Pfaffian expansion and the commutator
bound; it requires no choice of a global frame.

\textbf{Corollary 498.3 (conditional analytic vanishing).}
If there are smooth metrics $g_j$ for which
\[
 \int_M\left|\tfrac12g_j^{-1}\nabla g_j\right|_{g_j}^{2r}
       \,\mathrm d\operatorname{vol}_{g_j}\longrightarrow0,
\]
then $\chi(M)=0$. In fact a single metric making the right side of
\eqref{ra498:energy} smaller than one is enough, because $\chi(M)$ is
an integer.

\emph{Proof.}
Apply \eqref{ra498:energy} to each metric. The Euler integral is the same
topological integer for every metric, so its absolute value tends to zero.
The last assertion is the same inequality plus integrality.
\hfill$\square$

A naive scaling attempt fails exactly at the critical exponent. If
$\widehat g=c^2g$ for a positive constant $c$, then
$\widehat A=A$ as an endomorphism-valued one-form. Its endomorphism norm
is unchanged, its covector norm is multiplied by $c^{-1}$, and volume by
$c^{2r}$. Therefore
\[
 \int_M|\widehat A|_{\widehat g}^{2r}
       \,\mathrm d\operatorname{vol}_{\widehat g}
 =\int_M|A|_g^{2r}\,\mathrm d\operatorname{vol}_g.
\]
Making a metric uniformly large cannot achieve the hypothesis. Nor does
a partition of unity of local parallel metrics control it: differentiating
the weights produces nonzero $\nabla g$ in the overlaps. No estimate making
that critical energy arbitrarily small has been obtained here. The
criterion is sufficient and may be stronger than what the conjecture needs.

\paragraph{Why the missing ingredient must use more than flatness.}
Milnor's Theorem 2 [E4] supplies flat oriented real two-plane bundles with
nonzero Euler number: over a genus-two surface the Euler number one
satisfies $|1|<2$. Such a bundle is not that surface's tangent bundle,
whose Euler number is $-2$. It nevertheless disproves the general assertion
that every flat oriented bundle has zero real Euler class. Thus the metric
correction cannot automatically have an exact Euler form solely because
its starting bundle connection was flat.

For a tangent connection the extra torsion-free condition is an
integrability statement. In a local parallel frame $E_1,\ldots,E_n$,
\[
 T^\nabla(E_i,E_j)=\nabla_{E_i}E_j-\nabla_{E_j}E_i-[E_i,E_j]
                  =-[E_i,E_j].
\]
If torsion vanishes, the parallel frame fields commute and furnish local
affine coordinates. None of the algebra in Lemma 498.2 or the energy bound
uses this commutation. It is exactly the extra structure still available
to exploit, rather than a condition already spent in a nonexistent
curvature cancellation. Conversely, even using a local Hessian metric
does not force pointwise cancellation, as the explicit disk shows.

The remaining analytical target in this approach is therefore a global
construction using the affine-coordinate compatibility to produce either
the small-energy metrics of Corollary 498.3, or a global primitive for
$e(D)$, or a direct vanishing of its integral. No such construction is
provided. Local primitives on charts do not automatically glue: the
overlap corrections carry the Euler class one is trying to eliminate.
One cannot use Chern--Simons transgression from $\nabla$ as if the Euler
Pfaffian were an invariant polynomial for arbitrary general-linear
connections; that would repeat the initial error in a different notation.

\paragraph{Verification and outcome.}
The positive-dimensional exclusion prevents the point from being called
a counterexample. Finite covers preserve affine structures and have the
stated Euler scaling. The invariant-line proof does not assume compact
holonomy, and the compact-holonomy proof does not average over an infinite
measure. The metric identities were checked in parallel frames and then
interpreted tensorially; orthonormal frames are used only where the Euler
Pfaffian requires them. An exact rational matrix check in the local
research record verifies the Hessian, commutator, and nonzero curvature
at the origin. It is not a numerical experiment on the space of closed
affine manifolds.

\textbf{Outcome: UNRESOLVED.}
The work derives the metric-correction identity and its critical-energy
criterion, proves elementary holonomy and parity subcases, and applies the
known surface inequality. It also gives an explicit local obstruction to
a frequently tempting proof shortcut. The first unproved step is global
Euler-integral control in even dimensions beyond these subcases, using
torsion-free affine compatibility without assuming completeness or a
parallel volume form. No closed affine counterexample and no general
vanishing proof have been produced.

\subsection{Sources}
\begin{itemize}
\item[S1]Emmanuel Gnandi, Michel Nguiffo Boyom, and Stephane Puechmorel, Canonical foliations of statistical manifolds with statistical models, Remark 2.4.
\url{https://link.springer.com/article/10.1007/s41884-026-00193-8}.
\textit{Formulation source.}
\item[E1]Bruno Klingler, \emph{Chern's conjecture for special affine
manifolds}, Annals of Mathematics 186 (2017), no.~1, 69--95.
\url{https://annals.math.princeton.edu/2017/186-1/p02}.
\textit{Publisher theorem summary inspected 27 September 2026; the extra
parallel-volume hypothesis is essential to the cited statement.}
\item[E2]Jianquan Ge, \emph{Proof of Chern's conjecture on affine manifolds},
arXiv:2002.03105v2, withdrawn 27 July 2020.
\url{https://arxiv.org/abs/2002.03105v2}.
\textit{Withdrawal and author error notice inspected 27 September 2026;
not used as a theorem.}
\item[E3]Mihail Cocos, \emph{Proof of Chern conjecture for flat affine
manifolds}, arXiv:1504.04852v12, 6 December 2025.
\url{https://arxiv.org/abs/1504.04852v12}.
\textit{Record and abstract inspected 27 September 2026; general proof
claim not independently audited or used here.}
\item[E4]John Milnor, \emph{On the existence of a connection with curvature
zero}, Commentarii Mathematici Helvetici 32 (1958), 215--223, Theorems 1
and 2 and the surface corollary on the first page.
\url{https://doi.org/10.1007/BF02564579};
archival scan
\url{https://www.e-periodica.ch/cntmng?pid=com-001:1957:32::15}.
\textit{Indexed primary-paper theorem text inspected 27 September 2026;
direct archive retrieval encountered an access barrier.}
\item[E5]John W. Milnor and James D. Stasheff,
\emph{Characteristic Classes}, Princeton University Press, 1974,
Appendix C, Euler form and the distinction between orthogonal and general
linear connections.
\url{https://webhomes.maths.ed.ac.uk/~v1ranick/papers/milnstas.pdf}.
\textit{Indexed Appendix C passage inspected 27 September 2026;
the full scan exceeded the browser retrieval size limit.}
\end{itemize}
\end{document}
